\clearpage
\fancyhead[L]{\small CERES. PALLAS. JUNO. KOMET.}
\SourceLeaf{272}
\begin{center}
CERES, PALLAS, JUNO UND DER ZWEITE KOMET VON 1805.
\end{center}

\noindent\rule{\linewidth}{0.4pt}
\begin{center}
Astronomisches Jahrbuch für 1809. S. 137–140. Berlin 1806.
\end{center}
\noindent\rule{\linewidth}{0.4pt}

\begin{center}
GAUSS AN BODE.
\end{center}
\hfill Braunschweig, 14. März 1806.

\par\indent Gern teile ich Ew. --- einige Beobachtungen mit, die ich im vergangenen Winter angestellt habe und die Ihnen hoffentlich nicht unwillkommen sind. Es sind zwar nur wenige, doch das schlechte Wetter, über das alle meine astronomischen Korrespondenten klagen, hat auch hier den ganzen Winter über angehalten. Außerdem kann ich bei meinen dürftigen Hilfsmitteln die praktische Arbeit hier nur nebenbei betreiben. Zuerst folgen meine Beobachtungen der drei neuen Planeten.

\Needspace{13\baselineskip}
\begin{center}\emph{Ceres.}\end{center}
\begin{center}
\begingroup\fontsize{9}{11}\selectfont\setlength{\tabcolsep}{2pt}\renewcommand{\arraystretch}{1.1}
\begin{tabular}{lr|r|r}
\multicolumn{2}{c|}{Mittlere Zeit} & Rekt. & Dekl.\\ \hline
\textsuperscript{1}805. 24. Oktober & $11^{\mathrm u}\ 22^{\mathrm m}\ 32^{\mathrm s}$ & $109^\circ\ 28'\ 29''$ & $23^\circ\ 27'\ 0''$ Nord.\\
28. & $11\ 31\ 22$ & $110\ 7\ 7$ & $23\ 36\ 58$\\
14. November & $9\ 22\ 30$ & $111\ 44\ 56$ & $24\ 32\ 1$\\
18. & $9\ 0\ 37$ & $111\ 50\ 48$ & $24\ 48\ 47$\\
\end{tabular}\endgroup
\end{center}

\Needspace{11\baselineskip}
\begin{center}\emph{Pallas.}\end{center}
\begin{center}
\begingroup\fontsize{9}{11}\selectfont\setlength{\tabcolsep}{2pt}\renewcommand{\arraystretch}{1.1}
\begin{tabular}{lr|r|r}
1806. 14. Februar & $8^{\mathrm u}\ 11^{\mathrm m}\ 16^{\mathrm s}$ & $70^\circ\ 16'\ 31''$ & $19^\circ\ 59'\ 13''$ Süd.\\
16. & $7\ 32\ 28$ & $70\ 42\ 39$ & $19\ 20\ 44$\\
17. & $8\ 52\ 38$ & $70\ 56\ 44$ & $19\ 1\ 8$\\
20. & $7\ 49\ 35$ & $71\ 39\ 2$ & $18\ 5\ 0$\\
\end{tabular}\endgroup
\end{center}

\Needspace{8\baselineskip}
\begin{center}\emph{Juno.}\end{center}
\begin{center}
\begingroup\fontsize{9}{11}\selectfont\setlength{\tabcolsep}{2pt}\renewcommand{\arraystretch}{1.1}
\begin{tabular}{lr|r|r}
1806. 17. Februar & $9^{\mathrm u}\ 42^{\mathrm m}\ 0^{\mathrm s}$ & $173^\circ\ 46'\ 45''$ & $0^\circ\ 28'\ 32''$ Nord.\\
20. & $11\ 24\ 15$ & $173\ 15\ 53$ & $0\ 54\ 30$\\
\end{tabular}\endgroup
\end{center}

\par\indent Diesen Beobachtungen zufolge ergab meine Ephemeride für Ceres eine im Okt. 1805 um 4$\frac{1}{2}'$ und im Nov. um 5$\frac{1}{2}'$ zu große Rektaszension und eine um 1$'$ zu kleine Deklination. Für Pallas ergab sie eine um 27$\frac{1}{2}'$ zu große Rektaszension und eine um 5$'$ zu große Deklination (das heißt zu weit nördlich). Für Juno ergab sie eine um 1$\frac{1}{2}'$ zu große Rektaszension und eine um 1$'$ zu kleine Deklination (das heißt zu weit südlich). [Diese Beobachtungen der Ceres stehen auch in der Monatl. Corr., Band XIII, \href{https://www.digitale-sammlungen.de/en/view/bsb10538605?page=197}{S. 189}.]



\par\indent Der beträchtliche Unterschied bei Pallas ist zweifellos eine Folge der Störungen, die sich jetzt bemerkbar machen. Die über alle Erwartungen gute Übereinstimmung bei Juno beweist dagegen die Genauigkeit meiner letzten Elemente und meiner Beobachtung vom 20. Feb. 1805, nach der diese Elemente bestimmt wurden. Bisher kenne ich sonst niemanden, der Pallas oder Juno in diesem Jahr beobachtet oder gesehen hätte. Pallas ist sehr hell, so hell wie ein Stern 7. oder 7..8\textsuperscript{ter} Größe. Pallas muss sich auch bei voller Beleuchtung mit dem Passageninstrument und dem Quadranten gut beobachten lassen. Juno ist freilich lichtschwach. Bei meinen Beobachtungen hatte sie etwa 11\textsuperscript{ter} Größe und ist daher mit einem beleuchteten Fernrohr wohl kaum zu beobachten. Ceres ist wie immer gut sichtbar. Ich konnte meine Beobachtungen nur deshalb nicht fortsetzen, weil sie für meinen Beobachtungsplatz zu hoch kommt. Prof. \textsc{Harding} hat sie öfter gut am Mauerquadranten beobachtet. Gern werde ich Ihnen meine Ergebnisse und ihren Lauf für die nächste Erscheinung zusenden.

\SourceLeaf{273}

\par\indent Den zweiten Kometen des vergangenen Jahres habe ich nur am 8. Dezember beobachtet, und zwar zweimal. Er war auch mit bloßem Auge gut sichtbar, etwa so hell wie ein Stern 3. oder 4. Größe. Hier sind meine Beobachtungen.
\begin{center}
\begingroup\fontsize{9}{11}\selectfont\setlength{\tabcolsep}{2pt}\renewcommand{\arraystretch}{1.1}
\begin{tabular}{lr|r|r}
\multicolumn{2}{c|}{Mittlere Zeit} & Rektasz. & Dekl.\\ \hline
1805. 8. Dezember & $6^{\mathrm u}\ 55^{\mathrm m}\ 17^{\mathrm s}$ & $353^\circ\ 7'\ 40''$ & $23^\circ\ 36'\ 24''$ Südl.\\
& $7\ 58\ 3$ & $352\ 57\ 57$ & $24\ 6$\\
\end{tabular}\endgroup
\end{center}

\par\indent Die parabolischen Elemente, die ich aus dieser Beobachtung, einer Beobachtung von \textsc{Olbers} vom 2. Dez. und der im Moniteur veröffentlichten Beobachtung von \textsc{Bouvard} vom 16. Nov. berechnet habe, stehen im Januarheft der M. C. [\hyperlink{comet1805-2.1}{S. 266} dieses Bandes]. Inzwischen habe ich durch die Güte von Herrn \textsc{von Zach} noch die ganze Reihe der Originalbeobachtungen von \textsc{Thulis} erhalten. Außerdem ist durch die Zeitungen eine Meridianbeobachtung bekannt geworden, die am 8. Dez. in Greenwich gemacht wurde. Mit diesen zusätzlichen Beobachtungen habe ich die folgenden verbesserten Elemente gewonnen.

\Needspace{11\baselineskip}
\begin{center}
\begingroup\fontsize{8}{10}\selectfont\setlength{\tabcolsep}{2pt}\renewcommand{\arraystretch}{1.1}
\begin{tabular}{lr}
Zeit der Sonnennähe 1805. 31. Dez. . . . . . & $6^{\mathrm u}\ 54^{\mathrm m}\ 35^{\mathrm s}$ mittlere Seeberger Zeit\\
Länge der Sonnennähe . . . . . . . . . & $109^\circ\ 21'\ 50''5$\\
Länge des aufsteigenden Knotens . . . . . . & $250\ 33\ 34.9$\\
Neigung der Bahn . . . . . . . . . . . & $16\ 30\ 31.9$\\
Logarithmus des kleinsten Abstands . . . . & $9.9503300$\\
Bewegung . . . . . . . . . . . . . . & rechtläufig.\\
\end{tabular}\endgroup
\end{center}

\par\indent Hier folgt der Vergleich aller Beobachtungen mit diesen Elementen. Ich habe zwar noch ein paar weitere von \textsc{Thulis} erhalten, konnte sie aber nicht reduzieren, weil die zum Vergleich benutzten Sterne weder in den Sternverzeichnissen noch in der Hist. Cél. vorkommen. [Für \textsc{Gauss}, \textsc{Olbers} und \textsc{Thulis} ist jeweils die Zeit in Braunschweig, Bremen beziehungsweise Marseille angegeben.]
\begin{center}
\begingroup\fontsize{8}{10}\selectfont\setlength{\tabcolsep}{2pt}\renewcommand{\arraystretch}{1.1}
\begin{tabular}{r|r|r|r|r|r|l}
1805 & Mittlere Zeit & \shortstack{Berechnete\\Rektaszension} & \shortstack{Berechnete\\Deklination} & \shortstack{Differenz\\in Rekt.\\R---B} & \shortstack{Differenz\\in Dekl.\\R---B} & Beobachter\\ \hline
15. Nov. & $10^{\mathrm u}\ 28^{\mathrm m}\ 13^{\mathrm s}$ & $14^\circ\ 33'\ 0''$ & $39^\circ\ 18'\ 41''$ N. & $-2'\ 30''$ & $+\phantom{1'}\ 26''$ & Thulis\\
16 & $7\ 17\ 54$ & $14\ 12\ 7$ & $38\ 58\ 53$ & $-1\ 58$ & $+\phantom{1}\ 46$ & Thulis\\
19 & $7\ 11\ 17$ & $12\ 55\ 12$ & $37\ 34\ 27$ & $-\phantom{1}\ 4$ & $+1'\ 7$ & Thulis\\
20 & $7\ 3\ 15$ & $12\ 28\ 23$ & $36\ 59\ 41$ & $-\phantom{1}\ 23$ & $-\phantom{1}\ 32$ & Thulis\\
21 & $9\ 36\ 1$ & $11\ 57\ 5$ & $36\ 16\ 16$ & $+1\ 13$ & $-\phantom{1}\ 20$ & Thulis\\
29 & $6\ 39\ 44$ & $7\ 24\ 39$ & $26\ 47\ 49$ & $+1\ 32$ & $+\phantom{1}\ 1$ & Thulis\\
2. Dez. & $5\ 34\ 9$ & $4\ 48\ 43$ & $18\ 58\ 18$ & $+1\ 35$ & $-\phantom{1}\ 52$ & Olbers\\
2 & $9\ 39\ 0$ & $4\ 36\ 3$ & $18\ 20\ 28$ & $+1\ 47$ & $-\phantom{1}\ 3$ & Thulis\\
3 & $5\ 48\ 39$ & $3\ 39\ 51$ & $15\ 6\ 38$ & $+\phantom{1}\ 45$ & $-\phantom{1}\ 19$ & Olbers\\
3 & $6\ 30\ 37$ & $3\ 36\ 52$ & $14\ 55\ 11$ & $+1\ 10$ & $-\phantom{1}\ 26$ & Thulis\\
4 & $6\ 31\ 5$ & $2\ 16\ 36$ & $10\ 6\ 29$ N. & $+2\ 12$ & $-1\ 12$ & Thulis\\
6 & $7\ 16\ 16$ & $358\ 39\ 22$ & $3\ 39\ 32$ S. & $+\phantom{1}\ 22$ & $+\phantom{1}\ 48$ & Thulis\\
7 & $7\ 35\ 21$ & $356\ 9\ 21$ & $13\ 4\ 56$ & $\phantom{+1}\ 0$ & $+\phantom{1}\ 42$ & Thulis\\
8 & $5\ 27\ 22$ & $353\ 20\ 30$ & $22\ 56\ 24$ & $-1\ 10$ & $+2\ 44$ & Olbers\\
8 & $6\ 46\ 23$ & $353\ 8\ 33$ & $23\ 34\ 40$ & $-\phantom{1}\ 42$ & $+2\ 18$ & Olbers\\
8 & $6\ 55\ 17$ & $353\ 7\ 44$ & $23\ 35\ 57$ & $+\phantom{1}\ 4$ & $-\phantom{1}\ 47$ & Gauss\\
8 & & $353\ 6\ 21$ & $23\ 40\ 59$ & $-\phantom{1}\ 33$ & $-\phantom{1}\ 15$ & Maskelyne\\
8 & $6\ 59\ 39$ & $353\ 3\ 54$ & $23\ 47\ 34$ & $+\phantom{1}\ 9$ & $-\phantom{1}\ 33$ & Thulis\\
8 & $7\ 58\ 3$ & $352\ 57\ 34$ & $24\ 6\ 39$ & $-\phantom{1}\ 23$ & & Gauss\\
\end{tabular}\endgroup
\end{center}

\par\indent Die Elemente dieses Kometen liegen zwar sehr nahe bei denen des Kometen von 1772, doch ist der Unterschied, besonders beim kleinsten Abstand, immer noch zu groß. Eine eigene genaue Un%
\SourceLeaf{274}%
tersuchung der Beobachtungen des Kometen von 1772 hat mir außerdem gezeigt, dass man die beiden Sätze von Elementen einander nicht wesentlich näherbringen kann, ohne die Übereinstimmung mit den Beobachtungen zu beeinträchtigen. Der Komet von 1772 ist auch keinem Hauptplaneten so nahe gekommen, dass dessen Störung eine große Änderung seiner Elemente hätte verursachen können. [Diese Bemerkungen teilt \textsc{Gauss} nach \textsc{von Zach} am 21. Feb. 1806 mit (Monatl. Corr., Band XIII, \href{https://www.digitale-sammlungen.de/en/view/bsb10538605?page=320}{S. 310–313}, März 1806). Zugleich gibt er einige vorläufige Hinweise auf Arbeiten, deren Ergebnisse ausführlicher in dem weiter unten auf \hyperlink{native-275-lower}{S. 270} abgedruckten Brief vom 20. Mai 1806 stehen.]

\par\indent Ich arbeite jetzt an meiner Methode, die Planetenbahnen aus den Beobachtungen zu bestimmen, und werde hoffentlich bald daran denken können, sie öffentlich bekannt zu machen. Durch die vielen Verbesserungen, die ich seit 1801 einführen konnte, haben ihre verschiedenen Teile kaum noch Ähnlichkeit mit ihrer ursprünglichen Gestalt. Daher bereue ich es nicht, dass ich mich mit der Veröffentlichung nicht allzu sehr beeilt habe.

